% !PRINTED 376
% No folio or running head is printed on this chapter-opening leaf; printed page 376 is established by adjacent numbered leaves.
\weberpage{376}
\begin{center}
\gesperrt{Fünfzehnter Abschnitt.}

Geschlechter der quadratischen Formen.
\end{center}
\begin{center}\rule{0.2\linewidth}{0.4pt}\end{center}

\begin{center}
\textbf{\S\,103. \gesperrt{Darstellung von Zahlen durch quadratische Formen.}}
\end{center}

Eine natürliche Zahl $m$ heißt durch die primitive Form $(a,b,c)$ mit der Diskriminante $D$ \gesperrt{eigentlich darstellbar}, wenn es zwei relative Primzahlen $x$, $y$ gibt, die die Gleichung
\weq{1}{ax^2 + bxy + cy^2 = m}
erfüllen, und man kann eine mit $(a,b,c)$ äquivalente Form $(m,n,l)$ finden, deren Diskriminante
\weq{2}{D = n^2 - 4ml}
% [sic] Die Quelle verwendet trotz der zwei Klassen den Singular „diese Klasse“; der Wortlaut bleibt erhalten.
ist. Ist $m$ durch die Form $(a,b,c)$ darstellbar, so ist sie auch durch jede äquivalente Form, also durch jede Form der Klasse $A$, zu der $(a,b,c)$ gehört, darstellbar. Wir nennen dann $m$ durch die \gesperrt{Klasse $A$} darstellbar. Sind $A$, $A'$ zwei Klassen, $m$, $m'$ zwei zueinander teilerfremde Zahlen, die durch diese Klasse eigentlich darstellbar sind, so ist das Produkt $mm'$ durch die zusammengesetzte Klasse $AA'$ eigentlich darstellbar. Dies folgt aus den Kompositionen der Formen (12), (14), \S\,101.

Aus (2) ergab sich:
\weq{3}{n^2 \equiv D \ (\mathrm{mod}\ 4m).}

Ist $(m,n,l)$ imprimitiv, so kann $D$ keine Stammdiskriminante sein. Wir bezeichnen, wie bisher, den Stamm von $D$ mit $\Delta$ und setzen
\[
D = Q^2 \Delta,
\]
und der größte gemeinschaftliche Teiler von $m$, $n$, $l$ geht in $Q^2$ auf. Wir nehmen daher im Folgenden an, $m$ \gesperrt{sei relativ prim} zu $Q$, und sind dann sicher, daß die durch (2) \gesperrt{bestimmte Form $(m,n,l)$ primitiv} ist.
